\documentclass[10pt,a4paper]{amsart}
\usepackage[margin=19mm]{geometry}
\usepackage{amsmath,amssymb,mathtools}
\usepackage[scr=rsfs,cal=euler]{mathalfa}
\usepackage{xfrac}
\usepackage[unicode,hidelinks,bookmarks=false]{hyperref}
\newcommand{\ie}{i.e.\ }
\newcommand{\cf}{cf.\ }
\newcommand{\resp}{respectively\ }
\newcommand{\Subsection}{\subsection}
\providecommand{\nobreakdash}{\nobreak\hbox{-}}
\setlength{\emergencystretch}{2em}
\multlinegap0pt
\usepackage{bm}
\usepackage{bbm}
\usepackage{dsfont}
\usepackage{xspace}
\usepackage{cancel}
\usepackage{mathtools}
\usepackage{amsthm}
\makeatletter
\@ifundefined{remark}{\newtheorem{remark}{Remark}}{}
\@ifundefined{theorem}{\newtheorem{theorem}{Theorem}}{}
\@ifundefined{proposition}{\newtheorem{proposition}[theorem]{Proposition}}{}
\makeatother
\title[Topology and regularity for generalized ultradistribution al]{Topology and regularity for generalized ultradistribution algebras}
\author{Stevan Pilipovic, Dragana Risteski, Dimitris Scarpalezos, Milica Zigic}
\date{Source: arXiv:2411.07054v1 (CC BY 4.0)}
\begin{document}
\maketitle

\noindent\emph{Source and licence.} Topology and regularity for generalized ultradistribution algebras. Version arXiv:2411.07054v1 (CC BY 4.0), distributed under the Creative Commons Attribution 4.0 International licence, as recorded in the arXiv licence metadata for this exact version and shown on the abstract page https://arxiv.org/abs/2411.07054v1. Full rights notice: licenses/arxiv-2411.07054-CC-BY-4.0.txt.\par\smallskip

\noindent\textit{Editorial scope.} Counted unit: the Remark whose source label is \texttt{None} in the article (source0.tex lines 886--887, source label \texttt{None}), delivered together with the article's own complete original proof of it (source0.tex lines 888--943, one \texttt{proof} environment of 56 lines). No mathematics is written, repaired or re-derived: statement and proof are the article's own text line for line. Required context retained uncounted: thm 4.1 (lines 794--804); ops* (lines 878--882); p6 (lines 950--956). Every definition, display and label of those lines is retained unchanged. Omitted from this package and not redistributed: 1--793, 805--877, 883--885, 944--949, 957--1148. Nothing inside a retained excerpt is omitted or altered: every retained body is delivered exactly as the article's lines are written, so no omission receipt of any kind accompanies this package. Citations: the retained excerpts carry no citation command at all, so no bibliography entry of the article is transcribed and no reference is redistributed. Reference adaptation: none was needed and none was made; every reference inside the delivered unit resolves inside it, so no unresolved reference or citation is produced. Printed slips kept literal: no printed word, symbol, hypothesis or number of the counted statement or of its proof was altered. Layout and class: the article's own class is \texttt{sn-jnl} with packages including graphicx, multirow, amsmath, amssymb, amsfonts, amsthm, mathrsfs, appendix, xcolor, textcomp, manyfoot, booktabs, algorithm, algorithmicx; the wrapper is the lane's pinned \texttt{amsart} editorial wrapper at 10pt on A4, which restates the theorem-like environments the retained text needs and adds the article's own macro definitions that the retained text needs (none), with extra wrapper packages (bm, bbm, dsfont, xspace, cancel, mathtools). The delivered numbering is the unit's own. Assets: no retained span contains a figure environment, a graphics-inclusion command, a PDF-page inclusion, a table or a TikZ picture; no image file is copied into this package, the package declares no asset dependency and no image file is redistributed with it. Licence: version arXiv:2411.07054v1 of this article is distributed under the Creative Commons Attribution 4.0 International licence, as recorded in the arXiv licence metadata for this exact version and shown on the abstract page https://arxiv.org/abs/2411.07054v1. A full rights notice is delivered at licenses/arxiv-2411.07054-CC-BY-4.0.txt. Characters: every retained excerpt is pure ASCII, so the delivered engine, XeLaTeX with its default Latin Modern text font, reports no missing character.
\begin{theorem}\label{thm 4.1}
Let a generalized ultradistribution $[(f_\varepsilon)]\in\mathcal{G}^{(M_p)}(\Omega),$ respectively, $[(f_\varepsilon)]\in \mathcal{G}^{\{M_p\}}(\Omega)$ be such that $\mbox{supp }f_\varepsilon\subset K,$ $\varepsilon\in (0,1)$ and 
\begin{itemize}
\item in Beurling case: $\exists \tilde r>0,$ $\forall \phi\in \mathcal D^{M_p,\tilde r}(K),$ respectively,
\item in Roumieu case: $\exists (\tilde r_i)\in\mathfrak{R},$ $\forall \phi\in \mathcal D^{M_p,(\tilde r_i)}(K),$
\end{itemize}
\begin{align}\label{nulcon}
\left(\int_{\mathbb{R}^d}f_\varepsilon(x)\phi(x)\ dx\right)\in \mathcal{N}^{(M_p)}_\mathbb{C},\text{ respectively }\in \mathcal{N}^{\{M_p\}}_\mathbb{C}.
\end{align}
Then $f=0$ in $\mathcal{G}^{(M_p)}(\Omega),$ respectively, $f=0$ in $\mathcal{G}^{\{M_p\}}(\Omega).$
\end{theorem}

\begin{align}\label{ops*}
&\int f_\varepsilon(t)\kappa(t)dt\; e^{M(h/\varepsilon)}<C,\quad \varepsilon<\varepsilon_0,\quad %\nonumber\\
\text{respectively,}\\ \nonumber
&\int f_\varepsilon(t)\kappa(t)dt \;e^{N_{(h_i)}(1/\varepsilon)}<C,\quad \varepsilon<\varepsilon_0.
\end{align}

\begin{proposition}\label{p6}
In the Beurling case one can assume 
\begin{align*}
\forall h>0,\;\;\forall \kappa\in\mathcal{D}^{(M_p)}(K),\;\;&\exists \varepsilon_0\in (0,1),\;\;\exists C=C_{h,\kappa}>0,\\
&\int f_\varepsilon(t)\kappa(t)dt\;e^{M(h/\varepsilon)}\leq C,\quad \varepsilon<\varepsilon_0.
\end{align*}
\end{proposition}

\begin{remark} Condition \eqref{ops*} is stronger than $\left(\int f_\varepsilon(t)\kappa(t)dt\right) \in\mathcal N^\ast_\mathbb{C}$, for every $\kappa\in \mathcal{D}^\ast(\Omega).$ Actually, in the Beurling case we need to assume weaker condition than \eqref{ops*}. We will show this in Proposition \ref{p6} below. It is an open question whether we can in the Roumieu case assume a weaker condition than that in \eqref{ops*}.
\end{remark}

\begin{proof}
We assume that $\mbox{supp }f_\varepsilon\subset K,$ $ \varepsilon\in (0,1),$ 
where $K$ is a compact set of $\Omega.$ Our aim is to show that the conditions of Theorem \ref{thm 4.1} hold. We have 
\begin{align*}%\label{f51}
&\exists s>0,\; \forall h>0,\; \exists \tilde\varepsilon_0\in(0,1),\; \exists T_h=T>0,\;\; \left\| f_\varepsilon e^{-M(s/\varepsilon)}\right\|_{\mathcal{E}^{M_p,h}(K)} <T,\quad \varepsilon\in (0,\tilde\varepsilon_0),
\end{align*}
respectively,
\begin{align*}%\label{f52}
&\exists (s_i),\; \forall (h_i),\; \exists \tilde\varepsilon_0\in(0,1),\; \exists T_{(h_i)}=T>0,\;\; \left\| f_\varepsilon e^{-N_{(s_i)}(1/\varepsilon)}\right\|_{\mathcal{E}^{M_p,(h_i)}(K)} <T,\quad \varepsilon\in (0,\tilde\varepsilon_0).
\end{align*}
(We assume that $\varepsilon_0<\tilde\varepsilon_0$.) Condition \eqref{ops*} is equivalent with: there exists $\varepsilon_0$ such that for for every $h>0,$ respectively $(h_i)\in \mathfrak{R},$ 
\begin{align*} \label{nov11}
e^{M(h/\varepsilon)-M(s/\varepsilon)}\left\langle f_\varepsilon,\varphi\right\rangle=O(1),\;\; \varphi\in\mathcal D^{(M_p)}(\Omega),\text{ respectively, }\\
e^{N_{(h_i)}(1/\varepsilon)-N_{(s_i)}(1/\varepsilon)}\langle f_\varepsilon,\varphi\rangle=O(1),\;\; \varphi\in\mathcal D^{\{M_p\}}(\Omega),\;\varepsilon<\varepsilon_0.
\end{align*}
 This follows from the  properties of associated functions $M$ and $N_{(k_i)}.$

Let $\varepsilon<\varepsilon_0$. We note that for every $h>0$, respectively, $(h_i)$ there holds that
$e^{M(h/\varepsilon)-M(s/\varepsilon)}\langle f_\varepsilon, \psi \rangle$, respectively, $e^{N_{(h_i)}(1/\varepsilon)-N_{(s_i)}(1/\varepsilon)}\langle f_\varepsilon, \psi \rangle,$
is well defined for every $\psi\in \mathcal{D}^{M_p,h},$
respectively, for every $\psi\in \mathcal{D}^{M_p,(h_i)}.$

Next, fix $h>0$, respectively $(h_i)\in\mathfrak R.$  Since Banach spaces $\mathcal{D}^{M_p,h}(K)$ and $\mathcal{D}^{M_p,(h_i)}(K)$ are 
barreled and the set 
$\{e^{M(h/\varepsilon)-M(s/\varepsilon)}f_\varepsilon : \varepsilon<\varepsilon_0\}$, respectively 
$\{e^{N_{(h_i)}(1/\varepsilon)-N_{(s_i)}(1/\varepsilon)}f_\varepsilon:\varepsilon<\varepsilon_0\}$, is weakly bounded, it follows that it is strongly bounded. 
So, there exists open set $V_h=\{\varphi\in\mathcal{D}^{M_p,h}: \|\varphi\|_{K,h}<\delta\},$ respectively, open set $V_{(h_i)}=\{\varphi\in\mathcal{D}^{M_p,(h_i)}: \|\varphi\|_{K,(h_i)}<\delta\}$ such that 
\begin{align*}
\exists C>0,\;\forall \varphi\in V_h,\quad e^{M(k/\varepsilon)}\left\langle f_\varepsilon e^{-M(s/\varepsilon)},\varphi\right\rangle \leq C,\;\varepsilon<\varepsilon_0,
\end{align*}
respectively,
\begin{align*}
\exists C>0,\;\forall \varphi\in V_{(h_i)},\quad e^{N_{(h_i)}(1/\varepsilon)}\left\langle f_\varepsilon e^{-N_{(s_i)}(1/\varepsilon)},\varphi\right\rangle \leq C,\;\varepsilon<\varepsilon_0.
\end{align*}
We continue the proof in Beurling case. Let $(\varphi_n)_n$ be a sequence in $\mathcal{D}^{(M_p)}$ so that $\varphi_n\to \varphi$ in $\mathcal{D}^{M_p,h}$ ($\varphi\in \mathcal{D}^{M_p,h}$). There holds
\begin{align*}
e^{M(h/\varepsilon)}\left\langle f_\varepsilon e^{-M(s/\varepsilon)},\varphi_n-\varphi_m\right\rangle\leq C,\quad m,n>n_0, \;\varepsilon<\varepsilon_0.
\end{align*}
 This implies 
\begin{align*}
\left|\left\langle e^{M(h/\varepsilon)-M(s/\varepsilon)}f_\varepsilon,\varphi_n-\varphi\right\rangle\right|\leq C,\quad n>n_0,\; \varepsilon<\varepsilon_0.
\end{align*}
Thus for every $\varepsilon<\varepsilon_0$,
$
|\langle e^{M(h/\varepsilon)-M(s/\varepsilon)}f_\varepsilon,\varphi\rangle|\leq C.
$

If above instead of $M(h/\varepsilon),$ $M(s/\varepsilon)$ and $V_h$ one puts 
$N_{(h_i)}(1/\varepsilon),$ 
$N_{(s_i)}(1/\varepsilon)$ 
and $V_{(h_i)},$ one obtains that there exists $(h_i)\in \mathfrak{R}$ such that for every $\varphi\in \mathcal{D}^{M_p,(h_i)}$
\begin{align*}
\left|\left\langle e^{N_{(h_i)}(1/\varepsilon)-N_{(s_i)}(1/\varepsilon)}f_\varepsilon,\varphi\right\rangle\right|\leq C,\quad \varepsilon<\varepsilon_0.
\end{align*}
Thus, the conditions of the previous theorem holds and the assertion follows.
\end{proof}

\end{document}
